% Drop-in section; load section34_spectrum_figures.tex for its figure definitions.
\subsection{Bandwidth controls optimization access to high precision}
\label{sec:bandwidth-access}

The construction's nearly width-independent relative bandwidth
$\lambda=\gamma h$ requires $\gamma=\Theta(W)$.
Even though larger widths admit smaller $\lambda$ at comparable construction
error (Figure~2b), small slopes can delay readout optimization.
We quantify this in the dictionary of Section~3.1, replacing the direct
normalized least-squares solve with GD while keeping the features fixed.

\begin{theorem}[Slope-dependent output-error lower bound]
\label{thm:note-output}
For the fixed-center construction of Section~3.1 with common slope
$\gamma>0$ and nonzero sampled target, readout GD from zero with
$\eta=1/(2\mu_1)$ satisfies
\begin{equation}
 \frac{\|f-q_n\|_{2,Z}^{2}}{\|f\|_{2,Z}^{2}}
 \ge\sum_i p_i(\gamma)
 \left(1-\frac{\overline\rho_i(\gamma)}2\right)^{2n},\qquad n\ge0.
 \label{eq:note-output-bound}
\end{equation}
Here $\|\cdot\|_{2,Z}$ is the training-sample RMS norm,
$\mu_i$ are the descending readout-kernel eigenvalues, $p_i$ is the fraction
of target energy in eigenvector $i$, and
$\overline\rho_i(\gamma)\in[0,1]$ bounds $\mu_i/\mu_1$ from above.
The explicit gamma-dependent rate bounds and finite-geometry allowances
are given in Appendix~\ref{app:note-spectrum}.
\end{theorem}

\paragraph{Proof sketch.}
Gamma enters the feature correlations through $d\coth(\gamma d)$, where
$d$ is the separation between training points.
Controlling the finite-center corrections gives the relative-rate bounds;
a standard spectral argument then yields the output-error inequality.

\SectionSpectrumFigure
\FloatBarrier

Increasing $\lambda$ from $1/32$ to $1/4$ raises the thirtieth relative
eigenvalue by $1.4\times10^5$, versus only $1.3$ for the fourth
(Figure~\ref{fig:note-spectrum}); the thirtieth mode carries about 2.5\%
of target energy at both endpoints.

\SectionTrainingFigure
\FloatBarrier

At $\lambda=1/32$, a direct readout fit attains $6.2\times10^{-10}$ dense-grid
relative error, yet five million GD updates retain training error $0.210$
against the theorem's lower bound $0.2098$ (Figure~\ref{fig:note-training}A).
At $\lambda=1/4$, GD reaches $3.88\times10^{-4}$ over the same budget.
Across the four bandwidths, the bound remains at least 75\% of executed error
throughout training: small slopes delay accuracy already available in the
dictionary.

Over the same budget, joint Adam reaches median dense-grid relative error
$1.81\times10^{-3}$, versus $6.62\times10^{-7}$ on the supplied
$\lambda=1/4$ features, despite partial slope acquisition
(Figure~\ref{fig:note-training}B--C).
This motivates the next analysis: training must acquire useful slopes and
centers together. The construction supplies both the feature geometry and
the readout directly.
